\section{Czwarta metoda (\textbf{Metoda20})}
\subsection{Pomysł na metodę}
Metoda ta powstała na bazie idei integracji wielu modeli, podobnie jak w \textbf{Metoda17} i \textbf{Metoda18}, gdzie różne klasyfikatory były łączone w celu poprawy ogólnej wydajności predykcyjnej. Inspiracją były również metody stacking, takie jak \textbf{Metoda18}, gdzie predykcje modeli bazowych służyły jako cechy dla meta-modelu.

W \textbf{Metoda20} połączono LightGBM, CatBoost, Random Forest i Gradient Boosting. LightGBM i CatBoost to zaawansowane algorytmy boostingowe znane ze swojej szybkości i efektywności. Random Forest jest modelem opartym na zespołach losowo tworzonych predyktorów, a Gradient Boosting sekwencyjnie buduje modele korygujące błędy poprzedników. Podobnie jak w \textbf{Metoda18}, regresja logistyczna służy tutaj jako meta-model, integrując probabilistyczne predykcje tych czterech różnorodnych modeli bazowych. Celem jest wykorzystanie unikalnych mocnych stron każdego z tych algorytmów, aby osiągnąć wyższą dokładność i bardziej stabilne predykcje.

\subsection{Wejściowe dane}
Metoda \texttt{train\_run20} przyjmuje następujące parametry wejściowe:

\begin{enumerate}
	\def\labelenumi{\arabic{enumi}.}
	\item
	      \textbf{X\_train}: Macierz danych treningowych dla n próbek, m cech. Każdy wiersz reprezentuje próbkę, a każda kolumna cechę. Wartości w macierzy są zazwyczaj liczbowe.
	\item
	      \textbf{y\_train}: Wektor etykiet dla danych treningowych. Rozmiar: \textbf{n}, gdzie n - liczba próbek w danych treningowych.
	\item
	      \textbf{X\_test}: Macierz danych testowych, analogiczna do \textbf{X\_train}, używana do oceny modelu.
	\item
	      \textbf{y\_test}: Wektor etykiet testowych, odpowiadający danym testowym.
\end{enumerate}

\subsection{Wyjściowe dane}
Metoda zwraca następujące elementy:

\begin{enumerate}
	\def\labelenumi{\arabic{enumi}.}
	\item
	      \textbf{meta\_model}: Meta-model (regresja logistyczna) wytrenowany na predykcjach \textbf{LightGBM}, \textbf{CatBoost}, \textbf{Random Forest} i \textbf{Gradient Boosting}.
	\item
	      \textbf{meta\_features\_test}: Macierz meta-cech (predykcje modeli bazowych) dla danych testowych.
	\item
	      \textbf{y\_test}: Wektor etykiet testowych.
\end{enumerate}

\subsection{Algorytm}
\begin{enumerate}
	\def\labelenumi{\arabic{enumi}.}
	\item
	      Trenowanie modeli bazowych:
	      \begin{itemize}
		      \item Przeprowadza się trening \textbf{LGBMClassifier} na \textbf{X\_train} i \textbf{y\_train}.
		      \item Przeprowadza się trening \textbf{CatBoostClassifier} na \textbf{X\_train} i \textbf{y\_train}.
		      \item Przeprowadza się trening \textbf{RandomForestClassifier} na \textbf{X\_train} i \textbf{y\_train}.
		      \item Przeprowadza się trening \textbf{GradientBoostingClassifier} na \textbf{X\_train} i \textbf{y\_train}.
	      \end{itemize}
	\item
	      Generowanie predykcji dla danych treningowych:
	      \begin{itemize}
		      \item Uzyskuje się probabilistyczne predykcje klasy pozytywnej z każdego modelu bazowego na \textbf{X\_train}.
		      \item Tworzy się macierz meta-cech \textbf{meta\_features\_train}, gdzie każda kolumna zawiera predykcje jednego z modeli bazowych.
	      \end{itemize}
	\item
	      Trenowanie meta-modelu:
	      \begin{itemize}
		      \item Przeprowadza się trening \textbf{LogisticRegression} na \textbf{meta\_features\_train} i \textbf{y\_train}.
	      \end{itemize}
	\item
	      Generowanie predykcji dla danych testowych:
	      \begin{itemize}
		      \item Uzyskuje się probabilistyczne predykcje klasy pozytywnej z każdego wytrenowanego modelu bazowego na \textbf{X\_test}.
		      \item Tworzy się macierz meta-cech \textbf{meta\_features\_test}, analogicznie do treningowej.
	      \end{itemize}
	\item
	      Zwracanie wyników: Wracany jest wytrenowany \textbf{meta\_model}, \textbf{meta\_features\_test} i \textbf{y\_test}.
\end{enumerate}

\subsection{Kod}
\begin{lstlisting}[caption={Kod do 20 metody}, label={20Metoda}]
lgb_clf = LGBMClassifier(
n_estimators=100,
random_state=42
) # LightGBM
cat_clf = CatBoostClassifier(
iterations=100,
verbose=0,
random_state=42
) # CatBoost
rf_clf = RandomForestClassifier(
n_estimators=100,
random_state=42
) # Random Forest
gb_clf = GradientBoostingClassifier(
n_estimators=100,
random_state=42
) # Gradient Boosting

lgb_clf.fit(X_train, y_train)
cat_clf.fit(X_train, y_train)
rf_clf.fit(X_train, y_train)
gb_clf.fit(X_train, y_train)

lgb_preds_train = lgb_clf.predict_proba(X_train)[:, 1]
cat_preds_train = cat_clf.predict_proba(X_train)[:, 1]
rf_preds_train = rf_clf.predict_proba(X_train)[:, 1]
gb_preds_train = gb_clf.predict_proba(X_train)[:, 1]

meta_features_train = pd.DataFrame({
'lgb_preds': lgb_preds_train,
'cat_preds': cat_preds_train,
'rf_preds': rf_preds_train,
'gb_preds': gb_preds_train
})

meta_model = LogisticRegression()
meta_model.fit(meta_features_train, y_train)

lgb_preds_test = lgb_clf.predict_proba(X_test)[:, 1]
cat_preds_test = cat_clf.predict_proba(X_test)[:, 1]
rf_preds_test = rf_clf.predict_proba(X_test)[:, 1]
gb_preds_test = gb_clf.predict_proba(X_test)[:, 1]

meta_features_test = pd.DataFrame({
'lgb_preds': lgb_preds_test,
'cat_preds': cat_preds_test,
'rf_preds': rf_preds_test,
'gb_preds': gb_preds_test
})
return meta_model, meta_features_test, y_test
\end{lstlisting}

\subsection{Właściwości analityczne}
\begin{enumerate}
	\def\labelenumi{\arabic{enumi}.}
	\item
	      \textbf{Efektywność}: Łączenie predykcji wielu modeli (boostingowych i opartych na zespołach predyktorów) ma na celu zwiększenie dokładności i robustności predykcji.
	\item
	      \textbf{Wykorzystanie różnorodności modeli}: Zastosowanie różnych typów algorytmów pozwala na uwzględnienie różnych aspektów danych i wzorców.
	\item
	      \textbf{Meta-model}: Regresja logistyczna skutecznie integruje predykcje modeli bazowych, ucząc się, jak najlepiej je łączyć.
	\item
	      \textbf{Złożoność obliczeniowa}: Trening wielu złożonych modeli bazowych może być kosztowny obliczeniowo, ale predykcja jest stosunkowo szybka.
	\item
	      \textbf{Wydajność}: Oczekuje się, że kombinacja tych technik modelowania zapewni wysoką wydajność w zadaniu klasyfikacji.
\end{enumerate}